\chapter{Fault propagation as directed diffusion}
\label{ch:propagation}

The learning problem is only well posed if we say what a fault \emph{is}. We commit to a simple,
explicit generative model. Its purpose is twofold: it produces ground truth for training and
evaluation, and it is the physical prior that the network of Chapter~\ref{ch:symmetry} is asked to
invert. Readers with real incident data should replace this chapter with their data; the rest of
the document does not change.

\section{Injection}
A fault is a tuple $(v_0,t_0,s,T_r)$: origin node, onset step, peak severity, ramp length. Its
injection into the network is a ramp,
\begin{equation}
s(t)=s\cdot\min\!\Big(1,\ \frac{t-t_0+1}{T_r}\Big)\ind{t\ge t_0}.
\label{eq:injection}
\end{equation}
Long ramps are the \emph{silent} regime: the origin's own KPIs stay inside any static threshold for
tens of steps while the pattern across its neighbours is already distinctive.

\section{Nominal signal}
Each KPI follows a diurnal cycle plus coloured noise,
\begin{equation}
n_{vk}(t)=a_{vk}\sin\!\Big(\frac{2\pi t}{P}+\varphi_{vk}\Big)+\epsilon_{vk}(t),\qquad
\epsilon_{vk}(t)=r\,\epsilon_{vk}(t-1)+\eta_{vk}(t),\ \ \eta\sim\mathcal N(0,\sigma^2).
\label{eq:nominal}
\end{equation}
The AR(1) term is what makes per-node $z$-scores unreliable: consecutive excursions are correlated,
so a graph-free detector that fires on a short run of outliers has a high false-alarm rate
(Chapter~\ref{ch:experiments}).

\section{Observation model}
Degradation is a scalar latent field $d(t)\in\R^N_{\ge0}$. It enters the observed KPIs through a
type-specific effect direction $e_{\tau}\in\R^K$: a congested link shows latency and loss, a
sleeping cell shows throughput collapse, a signalling storm shows control-plane utilisation.
\begin{equation}
x_v(t)=n_v(t)+d_v(t)\,e_{\tau(v)}.
\label{eq:observation}
\end{equation}

\section{Propagation}
The field evolves as a damped, directed heat flow on the cause$\to$effect graph:
\begin{equation}
d(t+1)=\rho\,d(t)+\beta\,D_{\mathrm{in}}^{-1}A^{\top}d(t)+s(t)\,\mathbf{1}_{v_0},
\qquad (D_{\mathrm{in}})_{uu}=\max(1,\textstyle\sum_v A_{vu}).
\label{eq:propagation}
\end{equation}
$\rho$ is memory (a degraded router does not recover in one step), $\beta$ is coupling (how much of
a cause's degradation its effects inherit per step), and the in-degree normalisation says an
effect with many causes averages them rather than summing them.

\begin{proposition}[Steady state and attenuation]
With constant injection $s$ at an origin with no causes of its own, the origin reaches
$d_{v_0}^{\star}=s/(1-\rho)$, and a node $k$ hops downstream along a single chain reaches
$d^{\star}_{k}=\big(\tfrac{\beta}{1-\rho}\big)^{k}d^{\star}_{v_0}$.
Faults attenuate per hop iff
\begin{equation}
\frac{\beta}{1-\rho}<1.
\label{eq:stability}
\end{equation}
\end{proposition}
\begin{proof}
At steady state $d=\rho d+\beta P d+s\mathbf 1_{v_0}$ with $P=D_{\mathrm{in}}^{-1}A^\top$. For the
origin, $(Pd)_{v_0}=0$, giving the first claim. For a chain, $(Pd)_k=d_{k-1}$, so
$d_k(1-\rho)=\beta d_{k-1}$.
\end{proof}

We set $\rho=0.85$, $\beta=0.10$, so each hop keeps two thirds of its cause's degradation and the
blast radius of a core fault reaches the service layer with a recognisable but diminished signature.
The repository asserts \eqref{eq:stability} in its tests; violating it produces an exploding,
unphysical twin.

\begin{figure}[h]
\centering
\begin{tikzpicture}[x=0.09cm, y=1.1cm, >=Latex]
\draw[->, thick] (0,0) -- (160,0) node[below left, font=\scriptsize] {$t$};
\draw[->, thick] (0,0) -- (0,3.2) node[left, font=\scriptsize] {$d$};
\draw[dashed, ink!50] (50,0) -- (50,3) node[above, font=\scriptsize] {$t_0$};
\draw[dashed, rule] (0,2.5) -- (160,2.5) node[right, font=\scriptsize] {static threshold};
\draw[thick, tra] plot[smooth] coordinates {(0,0) (50,0) (60,0.6) (70,1.4) (80,2.1) (90,2.6) (100,2.9) (120,3.0) (160,3.0)} node[right, font=\scriptsize] {origin};
\draw[thick, ran] plot[smooth] coordinates {(0,0) (52,0) (62,0.3) (72,0.8) (82,1.3) (92,1.7) (102,1.9) (122,2.0) (160,2.0)} node[right, font=\scriptsize] {1 hop};
\draw[thick, svc] plot[smooth] coordinates {(0,0) (54,0) (64,0.15) (74,0.45) (84,0.8) (94,1.1) (104,1.25) (124,1.33) (160,1.33)} node[right, font=\scriptsize] {2 hops};
\draw[<->, thick, core] (68,2.65) -- (92,2.65) node[midway, above, font=\scriptsize] {lead time};
\end{tikzpicture}
\caption{The silent regime of \eqref{eq:injection}--\eqref{eq:propagation}. The origin crosses a
static threshold late; the joint pattern across hops is distinctive early. That gap is what the
graph buys.}
\end{figure}

\section{Why this is a diffusion and why that matters}
Equation \eqref{eq:propagation} is a discrete heat equation with a directed, degree-normalised
Laplacian $I-P$ and a source term. The forward map from injection history to the observed field is
linear and known. The learning problem is therefore not ``find a pattern'' but ``invert a known
operator from noisy, indirectly observed data'' (Chapter~\ref{ch:inverse}). That framing tells us
what the network must contain: the transposed relations, since $P^\top$ is the adjoint of the
propagation operator, and enough depth to span the support of the operator's impulse response.
